\section{模型验证与适用边界}\label{sec:validation}
\subsection{结果回放与独立反例}
验证分为三层。第一层是输入与对象一致性：核对原始文件摘要、编号唯一性、单位和资源库存，禁止用过期 Q3 输出驱动新输入下的分区。第二层是物理与调度回放：从货箱清单和访问顺序重新生成各架次的能耗、时间、交付完成时刻与 SOC，再检查资源占用区间和硬时限。第三层是模型边界反例，包括零距离、零载荷、最大载荷、电池恰好达到储备、DEM 边界与角点、覆盖内外和双向通信预算。

几何检查另采用独立矩形裁剪算法作为 DEM 像元穿越的交叉参照，而不是用同一函数给自己生成预期值。审计还通过篡改结果、替换输入、构造时间容差冲突等反例检查验证器能否发现错误。全仓库测试与正式数据全流程均重新运行，具体命令和日志随交付保存；没有把历史 PASS 作为本次通过证据。

需要区分“独立验证约束”与“独立物理模型”。回放独立检查结果表和资源逻辑，但能耗仍调用同一个官方物理引擎，因而不能排除所有共享公式误差。公式量纲检查与独立几何反例正是对这一盲区的补充；本文没有独立飞行实验可用。

\subsection{采样细化与最紧约束}
\paperfig{fig08_validation}{可行性与鲁棒性是不同层次的结论。左侧全部运输架次返航 SOC 超过 20\%，但最低仅 \MinSoc\%；右上为 \HardBoxes{} 个硬时限货箱的完成余量，最小 \HardSlack{} s；右下将通信步长由 1 s 减至 0.5 s，检查状态增加但最小链路余量下降。}{validation}
\papertable{sampling}{保持计划不变的通信采样细化检查；所有相位端点均纳入}{sampling}

1 s 设置检查 \StatesOne{} 个状态，0.5 s 设置检查 \StatesHalf{} 个状态，均未发现中断。最小余量由 \MarginOne{} dB 降至 \MarginHalf{} dB，说明更密检查确实找到了更紧状态，不能只比较中断次数。上述状态包括同一相位边界的重复记录，因此它们是检查记录数，不是独立统计样本量。

各相位内部采用不大于设定步长的等分采样，两个网格并不严格嵌套；不能据此次变化宣称一条普遍单调的“收敛曲线”。对连续通信的严格保证，还需要在遮挡状态变化处进行事件定位，或对整个时间区间给出可验证的链路余量界。目前这部分证明缺失，是 Q3 面对题设连续通信要求的明确不足。

\subsection{结果口径与可复现性}
表\ref{tab:summary}汇总不同层级的结果。Q1 时间为累计作业量；Q2 时间为运输任务完工；Q3 时间为运输与中继联合完工。能耗中 Q1/Q2 不含中继，Q3 含运输和中继，因而比较必须连同口径一起阅读。
\papertable{summary}{统一结果汇总；Q1 时间为累计量，Q2/Q3 时间为并行完工历时}{summary}

全部绘图脚本从结果 CSV 与派生诊断读取，不把最终结果数值写死在图中。表格和正文数值宏由同一生成脚本输出，官方结果模板也读取同一数据源。代码保持确定性排序；本次流程没有随机初始化、随机仿真或统计显著性检验，因而不虚构随机种子效应、置信区间或显著性。
\FloatBarrier
